\chapter{Saving of Results} 

\section{Node Results}
By default, the degree of freedom is saved to the universal file. But
sometimes one needs (additionally) the time derivative (e.g. in
acoustics, to calculate the pressure from the acoustic velocity
potential). Herefore, the command

\begin{verbatim}
<storeResults>
   <nodeResults type="velocity"/>
</storeResults>
\end{verbatim}

may be used. The value of the attribute {\bf type} is PDE dependent,
e.g. for mechanics, the following values are allowed:
\begin{itemize}
\item {\bf displacement}
\item {\bf velocity}    
\item {\bf acceleration}
\end{itemize}
The allowed types are mentioned at the individual PDE sections.

By default, the values on every node are stored. With the additional
attribute {\bf region} one can restrict the saved results to one or
more selected regions.
To store the results of the regions "bar" and "gap", one has to write
\begin{verbatim}
<storeResults>
   <nodeResults type="velocity" region="bar"/>
   <nodeResults type="velocity" region="gap"/>
</storeResults>
\end{verbatim}


\section{Element Results}
In general, saving of element results is, beside the different keyword
{\bf elemResults} totally similar to the procedure of saving nodal
results:
\begin{verbatim}
<storeResults>
   <elemResults type="stress"/>
</storeResults>
\end{verbatim}
Again, the values of {\bf type} depend on the regarded PDE. 


\section{Results at Special Points} 
In many applications we are only interested in the response of the
system at a single point. In order to reduce necessary disk space,
it is possible to store the nodal results only at a discrete
location. Therefore, the command {\bf nodeHistory} with the
attribute {\bf saveNodes="nodeName"} is provided. This command
reduces the amount of results stored tremendously, because no
universal-file is written.

In the config-file, the command is as follows

\begin{verbatim}
<storeResults>
   <nodeHistory type="velocity" saveNodes="ptName"/>
</storeResults>
\end{verbatim}

has to be used. The results are stored in the -- if necessary new
created -- directory {\em history} in a file with the name {\em
  conf-file-name.ptNr}.hist.

\begin{verbatim}
nsel,s,loc,x,len-1e-6,len+1e-6
nsel,r,loc,y,height
wsavnod,'ptName'
\end{verbatim}


\section{Results at Special Elements} 
Again, saving results of selected elements is similar to the procedure
of saving results on selected nodes. The commands name is {\bf
  elemHistory}, and the necessary keyword {\bf saveElems}:
\begin{verbatim}
   <storeResults>
      <elemHistory type="stress" saveElems="elName"/>
   </storeResults>
\end{verbatim}

